% Baseline vs persona merge -- SCHEMATIC, for intuition only. No data: answers
% are drawn as abstract text lines and the dots are illustrative, not measured.
% Caption it as a schematic.
% Preamble needs:  \usepackage{tikz}  \usetikzlibrary{calc}
% Include:  % Baseline vs persona merge -- SCHEMATIC, for intuition only. No data: answers
% are drawn as abstract text lines and the dots are illustrative, not measured.
% Caption it as a schematic.
% Preamble needs:  \usepackage{tikz}  \usetikzlibrary{calc}
% Include:  % Baseline vs persona merge -- SCHEMATIC, for intuition only. No data: answers
% are drawn as abstract text lines and the dots are illustrative, not measured.
% Caption it as a schematic.
% Preamble needs:  \usepackage{tikz}  \usetikzlibrary{calc}
% Include:  % Baseline vs persona merge -- SCHEMATIC, for intuition only. No data: answers
% are drawn as abstract text lines and the dots are illustrative, not measured.
% Caption it as a schematic.
% Preamble needs:  \usepackage{tikz}  \usetikzlibrary{calc}
% Include:  \input{figures/baseline_vs_ours}
%
% Colours match method_overview: blue = plain draft, aqua = persona A,
% orange = persona B (dataviz slots 1/3/2, validated all-pairs). Dots encode
% "represented" by FILL; colour only says which draft contributed.
\definecolor{cBlue}{HTML}{2A78D6}
\definecolor{cAqua}{HTML}{1BAF7A}
\definecolor{cOrange}{HTML}{EB6834}
\definecolor{cBase}{HTML}{6E6D68}
\definecolor{ink}{HTML}{0B0B0B}
\definecolor{ink2}{HTML}{52514E}
\definecolor{rule}{HTML}{C9C8C3}
\definecolor{line}{HTML}{BDBCB6}

\begin{tikzpicture}[font=\sffamily\footnotesize, x=1cm, y=1cm]

% "text" = a stack of grey bars; #1 x, #2 y-top, #3 width, #4 number of lines
\newcommand{\txt}[4]{%
  \foreach \k in {1,...,#4} {
    \pgfmathsetmacro{\w}{ifthenelse(\k==#4, 0.62, ifthenelse(mod(\k,3)==0, 0.88, 1))}
    \fill[line, rounded corners=0.6pt] (#1,#2-0.17*\k+0.04) rectangle ++(#3*\w,0.07);
  }}
% cluster dots: every dot is outlined (= a cluster); represented ones are filled
\newcommand{\clrow}[2]{% #1 x of first dot, #2 y
  \foreach \i in {0,...,5} {\draw[ink2, line width=0.8pt] (#1+0.4*\i,#2) circle (0.13);}}
\newcommand{\clfill}[3]{% #1 x of first dot, #2 y, #3 list of index/colour
  \foreach \i/\c in {#3} {\fill[\c] (#1+0.4*\i,#2) circle (0.13);}}

% question
\node[draw=ink, rounded corners=3pt, inner sep=5pt, font=\sffamily\small]
  at (7.0,0.45) {\textbf{Question}\quad\emph{Do you believe in God?}};

% ======================================================= baseline (left)
\node[font=\sffamily\small\bfseries, text=cBase, anchor=west] at (0,-0.35)
  {Single answer};
\draw[cBase, line width=0.8pt, rounded corners=3pt, fill=cBase!5]
  (0,-0.6) rectangle (6.2,-2.95);
\txt{0.3}{-0.75}{5.6}{12}
\node[font=\sffamily\scriptsize, text=ink2, anchor=north west, align=left]
  at (0.05,-3.0) {one balanced answer, pitched at the middle};

\node[font=\sffamily\scriptsize, text=ink2, anchor=west] at (0,-3.65)
  {viewpoint clusters represented};
\clrow{3.8}{-3.65}
\clfill{3.8}{-3.65}{2/cBase, 3/cBase}

% ========================================================== ours (right)
\node[font=\sffamily\small\bfseries, text=cBlue, anchor=west] at (7.8,-0.35)
  {Persona merge (ours)};
\draw[rule, line width=0.8pt, rounded corners=3pt, fill=white]
  (7.8,-0.6) rectangle (14.0,-2.95);
% three paragraphs, each with a coloured rule on the left
\foreach \y/\c/\lab/\n in {-0.72/cBlue/plain/3, -1.48/cAqua/group A's view/3,
                           -2.24/cOrange/group B's view/3} {
  \fill[\c] (7.95,\y-0.02) rectangle (8.03,\y-0.6);
  \node[font=\sffamily\scriptsize\bfseries, text=\c, anchor=north west]
    at (8.12,\y+0.03) {\lab};
  \txt{10.25}{\y+0.02}{3.55}{\n}
}
\node[font=\sffamily\scriptsize, text=ink2, anchor=north west, align=left]
  at (7.85,-3.0) {one answer, a paragraph per view};

\node[font=\sffamily\scriptsize, text=ink2, anchor=west] at (7.8,-3.65)
  {viewpoint clusters represented};
\clrow{11.6}{-3.65}
\clfill{11.6}{-3.65}{0/cAqua, 2/cBlue, 3/cBlue, 4/cOrange}
\end{tikzpicture}

%
% Colours match method_overview: blue = plain draft, aqua = persona A,
% orange = persona B (dataviz slots 1/3/2, validated all-pairs). Dots encode
% "represented" by FILL; colour only says which draft contributed.
\definecolor{cBlue}{HTML}{2A78D6}
\definecolor{cAqua}{HTML}{1BAF7A}
\definecolor{cOrange}{HTML}{EB6834}
\definecolor{cBase}{HTML}{6E6D68}
\definecolor{ink}{HTML}{0B0B0B}
\definecolor{ink2}{HTML}{52514E}
\definecolor{rule}{HTML}{C9C8C3}
\definecolor{line}{HTML}{BDBCB6}

\begin{tikzpicture}[font=\sffamily\footnotesize, x=1cm, y=1cm]

% "text" = a stack of grey bars; #1 x, #2 y-top, #3 width, #4 number of lines
\newcommand{\txt}[4]{%
  \foreach \k in {1,...,#4} {
    \pgfmathsetmacro{\w}{ifthenelse(\k==#4, 0.62, ifthenelse(mod(\k,3)==0, 0.88, 1))}
    \fill[line, rounded corners=0.6pt] (#1,#2-0.17*\k+0.04) rectangle ++(#3*\w,0.07);
  }}
% cluster dots: every dot is outlined (= a cluster); represented ones are filled
\newcommand{\clrow}[2]{% #1 x of first dot, #2 y
  \foreach \i in {0,...,5} {\draw[ink2, line width=0.8pt] (#1+0.4*\i,#2) circle (0.13);}}
\newcommand{\clfill}[3]{% #1 x of first dot, #2 y, #3 list of index/colour
  \foreach \i/\c in {#3} {\fill[\c] (#1+0.4*\i,#2) circle (0.13);}}

% question
\node[draw=ink, rounded corners=3pt, inner sep=5pt, font=\sffamily\small]
  at (7.0,0.45) {\textbf{Question}\quad\emph{Do you believe in God?}};

% ======================================================= baseline (left)
\node[font=\sffamily\small\bfseries, text=cBase, anchor=west] at (0,-0.35)
  {Single answer};
\draw[cBase, line width=0.8pt, rounded corners=3pt, fill=cBase!5]
  (0,-0.6) rectangle (6.2,-2.95);
\txt{0.3}{-0.75}{5.6}{12}
\node[font=\sffamily\scriptsize, text=ink2, anchor=north west, align=left]
  at (0.05,-3.0) {one balanced answer, pitched at the middle};

\node[font=\sffamily\scriptsize, text=ink2, anchor=west] at (0,-3.65)
  {viewpoint clusters represented};
\clrow{3.8}{-3.65}
\clfill{3.8}{-3.65}{2/cBase, 3/cBase}

% ========================================================== ours (right)
\node[font=\sffamily\small\bfseries, text=cBlue, anchor=west] at (7.8,-0.35)
  {Persona merge (ours)};
\draw[rule, line width=0.8pt, rounded corners=3pt, fill=white]
  (7.8,-0.6) rectangle (14.0,-2.95);
% three paragraphs, each with a coloured rule on the left
\foreach \y/\c/\lab/\n in {-0.72/cBlue/plain/3, -1.48/cAqua/group A's view/3,
                           -2.24/cOrange/group B's view/3} {
  \fill[\c] (7.95,\y-0.02) rectangle (8.03,\y-0.6);
  \node[font=\sffamily\scriptsize\bfseries, text=\c, anchor=north west]
    at (8.12,\y+0.03) {\lab};
  \txt{10.25}{\y+0.02}{3.55}{\n}
}
\node[font=\sffamily\scriptsize, text=ink2, anchor=north west, align=left]
  at (7.85,-3.0) {one answer, a paragraph per view};

\node[font=\sffamily\scriptsize, text=ink2, anchor=west] at (7.8,-3.65)
  {viewpoint clusters represented};
\clrow{11.6}{-3.65}
\clfill{11.6}{-3.65}{0/cAqua, 2/cBlue, 3/cBlue, 4/cOrange}
\end{tikzpicture}

%
% Colours match method_overview: blue = plain draft, aqua = persona A,
% orange = persona B (dataviz slots 1/3/2, validated all-pairs). Dots encode
% "represented" by FILL; colour only says which draft contributed.
\definecolor{cBlue}{HTML}{2A78D6}
\definecolor{cAqua}{HTML}{1BAF7A}
\definecolor{cOrange}{HTML}{EB6834}
\definecolor{cBase}{HTML}{6E6D68}
\definecolor{ink}{HTML}{0B0B0B}
\definecolor{ink2}{HTML}{52514E}
\definecolor{rule}{HTML}{C9C8C3}
\definecolor{line}{HTML}{BDBCB6}

\begin{tikzpicture}[font=\sffamily\footnotesize, x=1cm, y=1cm]

% "text" = a stack of grey bars; #1 x, #2 y-top, #3 width, #4 number of lines
\newcommand{\txt}[4]{%
  \foreach \k in {1,...,#4} {
    \pgfmathsetmacro{\w}{ifthenelse(\k==#4, 0.62, ifthenelse(mod(\k,3)==0, 0.88, 1))}
    \fill[line, rounded corners=0.6pt] (#1,#2-0.17*\k+0.04) rectangle ++(#3*\w,0.07);
  }}
% cluster dots: every dot is outlined (= a cluster); represented ones are filled
\newcommand{\clrow}[2]{% #1 x of first dot, #2 y
  \foreach \i in {0,...,5} {\draw[ink2, line width=0.8pt] (#1+0.4*\i,#2) circle (0.13);}}
\newcommand{\clfill}[3]{% #1 x of first dot, #2 y, #3 list of index/colour
  \foreach \i/\c in {#3} {\fill[\c] (#1+0.4*\i,#2) circle (0.13);}}

% question
\node[draw=ink, rounded corners=3pt, inner sep=5pt, font=\sffamily\small]
  at (7.0,0.45) {\textbf{Question}\quad\emph{Do you believe in God?}};

% ======================================================= baseline (left)
\node[font=\sffamily\small\bfseries, text=cBase, anchor=west] at (0,-0.35)
  {Single answer};
\draw[cBase, line width=0.8pt, rounded corners=3pt, fill=cBase!5]
  (0,-0.6) rectangle (6.2,-2.95);
\txt{0.3}{-0.75}{5.6}{12}
\node[font=\sffamily\scriptsize, text=ink2, anchor=north west, align=left]
  at (0.05,-3.0) {one balanced answer, pitched at the middle};

\node[font=\sffamily\scriptsize, text=ink2, anchor=west] at (0,-3.65)
  {viewpoint clusters represented};
\clrow{3.8}{-3.65}
\clfill{3.8}{-3.65}{2/cBase, 3/cBase}

% ========================================================== ours (right)
\node[font=\sffamily\small\bfseries, text=cBlue, anchor=west] at (7.8,-0.35)
  {Persona merge (ours)};
\draw[rule, line width=0.8pt, rounded corners=3pt, fill=white]
  (7.8,-0.6) rectangle (14.0,-2.95);
% three paragraphs, each with a coloured rule on the left
\foreach \y/\c/\lab/\n in {-0.72/cBlue/plain/3, -1.48/cAqua/group A's view/3,
                           -2.24/cOrange/group B's view/3} {
  \fill[\c] (7.95,\y-0.02) rectangle (8.03,\y-0.6);
  \node[font=\sffamily\scriptsize\bfseries, text=\c, anchor=north west]
    at (8.12,\y+0.03) {\lab};
  \txt{10.25}{\y+0.02}{3.55}{\n}
}
\node[font=\sffamily\scriptsize, text=ink2, anchor=north west, align=left]
  at (7.85,-3.0) {one answer, a paragraph per view};

\node[font=\sffamily\scriptsize, text=ink2, anchor=west] at (7.8,-3.65)
  {viewpoint clusters represented};
\clrow{11.6}{-3.65}
\clfill{11.6}{-3.65}{0/cAqua, 2/cBlue, 3/cBlue, 4/cOrange}
\end{tikzpicture}

%
% Colours match method_overview: blue = plain draft, aqua = persona A,
% orange = persona B (dataviz slots 1/3/2, validated all-pairs). Dots encode
% "represented" by FILL; colour only says which draft contributed.
\definecolor{cBlue}{HTML}{2A78D6}
\definecolor{cAqua}{HTML}{1BAF7A}
\definecolor{cOrange}{HTML}{EB6834}
\definecolor{cBase}{HTML}{6E6D68}
\definecolor{ink}{HTML}{0B0B0B}
\definecolor{ink2}{HTML}{52514E}
\definecolor{rule}{HTML}{C9C8C3}
\definecolor{line}{HTML}{BDBCB6}

\begin{tikzpicture}[font=\sffamily\footnotesize, x=1cm, y=1cm]

% "text" = a stack of grey bars; #1 x, #2 y-top, #3 width, #4 number of lines
\newcommand{\txt}[4]{%
  \foreach \k in {1,...,#4} {
    \pgfmathsetmacro{\w}{ifthenelse(\k==#4, 0.62, ifthenelse(mod(\k,3)==0, 0.88, 1))}
    \fill[line, rounded corners=0.6pt] (#1,#2-0.17*\k+0.04) rectangle ++(#3*\w,0.07);
  }}
% cluster dots: every dot is outlined (= a cluster); represented ones are filled
\newcommand{\clrow}[2]{% #1 x of first dot, #2 y
  \foreach \i in {0,...,5} {\draw[ink2, line width=0.8pt] (#1+0.4*\i,#2) circle (0.13);}}
\newcommand{\clfill}[3]{% #1 x of first dot, #2 y, #3 list of index/colour
  \foreach \i/\c in {#3} {\fill[\c] (#1+0.4*\i,#2) circle (0.13);}}

% question
\node[draw=ink, rounded corners=3pt, inner sep=5pt, font=\sffamily\small]
  at (7.0,0.45) {\textbf{Question}\quad\emph{Do you believe in God?}};

% ======================================================= baseline (left)
\node[font=\sffamily\small\bfseries, text=cBase, anchor=west] at (0,-0.35)
  {Single answer};
\draw[cBase, line width=0.8pt, rounded corners=3pt, fill=cBase!5]
  (0,-0.6) rectangle (6.2,-2.95);
\txt{0.3}{-0.75}{5.6}{12}
\node[font=\sffamily\scriptsize, text=ink2, anchor=north west, align=left]
  at (0.05,-3.0) {one balanced answer, pitched at the middle};

\node[font=\sffamily\scriptsize, text=ink2, anchor=west] at (0,-3.65)
  {viewpoint clusters represented};
\clrow{3.8}{-3.65}
\clfill{3.8}{-3.65}{2/cBase, 3/cBase}

% ========================================================== ours (right)
\node[font=\sffamily\small\bfseries, text=cBlue, anchor=west] at (7.8,-0.35)
  {Persona merge (ours)};
\draw[rule, line width=0.8pt, rounded corners=3pt, fill=white]
  (7.8,-0.6) rectangle (14.0,-2.95);
% three paragraphs, each with a coloured rule on the left
\foreach \y/\c/\lab/\n in {-0.72/cBlue/plain/3, -1.48/cAqua/group A's view/3,
                           -2.24/cOrange/group B's view/3} {
  \fill[\c] (7.95,\y-0.02) rectangle (8.03,\y-0.6);
  \node[font=\sffamily\scriptsize\bfseries, text=\c, anchor=north west]
    at (8.12,\y+0.03) {\lab};
  \txt{10.25}{\y+0.02}{3.55}{\n}
}
\node[font=\sffamily\scriptsize, text=ink2, anchor=north west, align=left]
  at (7.85,-3.0) {one answer, a paragraph per view};

\node[font=\sffamily\scriptsize, text=ink2, anchor=west] at (7.8,-3.65)
  {viewpoint clusters represented};
\clrow{11.6}{-3.65}
\clfill{11.6}{-3.65}{0/cAqua, 2/cBlue, 3/cBlue, 4/cOrange}
\end{tikzpicture}
